\section{推理的局限性}

\subsection{LLM 容易被无关上下文干扰}

\begin{itemize}
    \item 在数学题中插入无关信息（如 ``Mario's moustache costs \$10''），LLM 会被误导。
    \item 添加提示 ``ignore irrelevant context'' 可以部分缓解，但当无关内容量增大时效果有限。
    \item 即使是简单的无关句子（``the sky is blue''）大量堆积，也会导致显著的性能下降。
\end{itemize}

\begin{warningbox}{上下文干扰}
LLM 作为概率模型，会将所有输入 token 纳入注意力计算。无关信息不仅浪费上下文窗口，还会实质性地干扰推理过程。这在实际的 Agent 应用中需要特别注意——检索增强（RAG）引入的噪声文档可能反而损害推理质量。
\end{warningbox}

\subsection{LLM 无法自我纠错推理}

Denny Zhou 团队的实验表明：

\begin{itemize}
    \item 让 LLM 审查并修正自己的答案时，它可能纠正错误答案，但也会把正确答案改错。
    \item 在 GSM8K、CommonsenseQA 等基准上，self-correction 方法\textbf{没有带来任何净提升}，反而使性能变差。
    \item 文献中报告的 self-correction 提升往往使用了 Oracle（即只在答案错误时才要求模型修正）——但模型本身无法判断自己是否正确。
\end{itemize}

\begin{importantbox}{外部反馈是自我纠错的前提}
LLM 的自我纠错需要外部 Oracle 反馈（如代码任务中的 unit test）。Self-Debug 工作通过单元测试自然地提供了这种 Oracle。纯粹依赖模型自身判断来纠错是不可靠的。
\end{importantbox}

\subsection{Multi-Agent Debate 并不优于 Self-Consistency}

\begin{itemize}
    \item 多个 Agent 互相辩论并达成共识（multi-agent debate），总共生成 $n$ 个回答。
    \item 但简单地对 $n$ 个独立采样使用 Self-Consistency 投票，效果始终优于 multi-agent debate。
    \item 辩论过程引入的信息交互并未带来额外收益。
\end{itemize}

\subsection{前提顺序（Premise Order）影响推理}

\begin{itemize}
    \item \textbf{GSM8K 实验}：仅将问题中的句子重新排列（不改变语义），准确率下降约 10 个百分点。
    \item \textbf{逻辑推理实验}：使用随机符号的纯逻辑规则推理，仅打乱相关规则的顺序，所有 Frontier 模型准确率下降 30+ 个百分点。
    \item \textbf{根本原因}：LLM 只能顺序处理输入，无法像人类一样在前提之间自由跳转和回溯。
\end{itemize}

\begin{warningbox}{前提顺序效应}
LLM 的推理严重依赖前提的呈现顺序。在设计 Agent 系统时，输入信息的组织顺序需要与推理所需的顺序对齐，否则会导致显著的性能损失。
\end{warningbox}

\subsection{本章小结}

当前 LLM 推理存在三大局限：容易被无关上下文干扰、无法可靠自我纠错、对前提顺序敏感。这些局限性为 Agent 系统的设计提供了重要指导——需要通过外部工具和精心设计的工作流来弥补这些短板。

